% --- LaTeX Essay Template - S. Venkatraman ---

% --- Set document class and font size ---

\documentclass[letterpaper, 11pt]{article}

% --- Package imports ---

% lipsum is just for generating placeholder text and can be removed
\usepackage{hyperref, lipsum} 

% --- Fonts (Set to Palatino or Times New Roman if desired) ---

% \usepackage{newpxtext, newpxmath}
% \usepackage{times}

% --- Page layout settiings ---

% Set page margins
\usepackage[left=1.2in, right=1.2in, top=.9in, bottom=.9in]{geometry}

% Anchor footnotes to the bottom of the page
\usepackage[bottom]{footmisc}
\usepackage{setspace}
\setstretch{1.15}
% Set line spacing
\renewcommand{\baselinestretch}{1.2}

% Use this for a one-off '*' footnote (e.g., a URL in a heading)
\newcommand{\starfootnote}[1]{%
\begingroup
\renewcommand{\thefootnote}{\fnsymbol{footnote}}%
\footnote{#1}%
\endgroup
\addtocounter{footnote}{-1}%
}

% Set spacing between paragraphs
\setlength{\parskip}{2mm}

% --- Page formatting settings ---



% --- Essay title and author ---
\title{\vspace{-1.5cm} Summary of \textit{Veridical Data Science}, Chapters 10 \footnote{\url{https://vdsbook.com/}}}

\date{\today}

% --- Main text ---
% --- Document starts here ---

\begin{document}
% Set link colors for labeled items and URLs (blue) 
\hypersetup{colorlinks=true, linkcolor=blue, urlcolor=blue}
\maketitle

\section*{Chapter 10: Extending the Least Squares Algorithm}

Chapter 10 introduces feature engineering and regularization in predictive modeling.
It explains that raw data often require preprocessing, including encoding categorical variables (one-hot encoding) and scaling predictors.

The chapter discusses overfitting as capturing noise rather than signal.
Ridge (L2) and lasso (L1) regularization could help us avoid overfitting by adding a penalty to the loss function that discourages large coefficients.
Ridge shrinks coefficients towards zero, while lasso can set some coefficients exactly to zero, performing feature selection.

Feature selection and preprocessing decisions must be made using training data to avoid data leakage and ensure valid model evaluation.
And cross-validation is essential for tuning hyperparameters and assessing generalization performance.

% --- Document ends here ---
\end{document}
